% =============================================================================
\chapter{Rosenbaum's \texorpdfstring{$\Gamma$}{Gamma}, and Table 1 on real structure (M1)}\label{ch:gamma}

Chapters~\ref{ch:frame} to~\ref{ch:pest} audited a certificate. This chapter is about what the
paper is for, and about the first measurement that answers that question on real structures.

\begin{ideia}[The paper in two sentences]
Rosenbaum did not test what the data cannot test: he asked how strong hidden confounding would
have to be to overturn the conclusion. We ask the same question of background knowledge, and the
unit of violation is the number of the analyst's orientation claims that would have to be withdrawn.
\end{ideia}

\begin{antes}
\antesq[deriva-enquadramento]{The first measurement-paper draft of 12 September had correct
numbers. What was wrong with it?}
\antesq[gamma-de-rosenbaum]{What does $\Gamma$ bound, in what design, and what does $\Gamma = 1$
mean?}
\end{antes}

\section{12 September: the subject was wrong}

On 12 September, with Table~4 of the ledger just out, I wrote a full draft of the measurement paper,
\emph{Claim Radius: Certifying Covariate Adjustment Sets Against Misspecified Background Knowledge}.
Its numbers came from the result files, and its centre was the supplier ledger: knowledge suppliers
from 7B to 72B, scrambled names, a scale axis as a table of its own. That was the wrong paper. The
evaluation protocol had been built first, and its most elaborate arm had become the subject. What we
set out to write is an uncertainty measure for background knowledge, the analogue of Rosenbaum's
$\Gamma$ for orientation claims. Elicited knowledge is one of four knowledge processes in the
evaluation (with controlled reversals, chance orientations and truthful controls, the last asserting nothing false), never the title,
the abstract's lead or Table~1. Your related-work paragraph (\texttt{sections/03\_related\_work.tex})
already has the sentence: the radius asks the structural analogue of how strong a confounder must be.
The change is that this sentence becomes the spine of the paper.

\textbf{What Rosenbaum does.} In a matched study each pair joins a treated and a control unit with
the same observed covariates; with perfect matching either member is equally likely to be the treated
one. An unmeasured covariate breaks this. Rosenbaum does not estimate it: he lets the two members
differ in their odds of treatment by at most a factor $\Gamma$, redoes the inference in the worst case
the model allows, and reports the $\Gamma$ at which the conclusion is lost.

\begin{exemplo}[Illustrative: the numbers are constructed for this page, not measured in any study]
Take 40 discordant pairs with a binary outcome; in each, exactly one member had the event, and in
30 of them it was the treated one. With no treatment effect and $\Gamma = 1$, the count $T$ of
pairs where the treated member had the event is $\mathrm{Bin}(40, 1/2)$ and the one-sided p-value is
$P(T \ge 30) \approx 0.001$. With $\Gamma > 1$ the chance that the event member is the treated one
may be as large as $p^{+} = \Gamma/(1+\Gamma)$, and the worst-case p-value is the upper tail of
$\mathrm{Bin}(40, p^{+})$. For $\Gamma = 1, 1.25, 1.5, 1.6, 1.75$ it is $0.001, 0.009, 0.035, 0.053,
0.089$ (exact binomial tails, computed for this page). Significance at 5\,\% is lost near
$\Gamma^{*} \approx 1.58$, and the report reads: the conclusion survives a hidden confounder that
multiplies the within-pair odds of treatment by up to 1.58, and falls above that.
\end{exemplo}

The example is the sign test, the simplest case of this model.

\begin{definicao}[Sensitivity model and sensitivity value]
Let $\pi_{1}, \pi_{2}$ be the treatment probabilities of two members of a pair with the same observed
covariates. The model with parameter $\Gamma \ge 1$ admits every assignment with
\[
\frac{1}{\Gamma} \;\le\; \frac{\pi_{1}(1-\pi_{2})}{\pi_{2}(1-\pi_{1})} \;\le\; \Gamma .
\]
Given that exactly one member is treated, the probability that it is the first lies in
$[1/(1+\Gamma),\ \Gamma/(1+\Gamma)]$, and $\Gamma = 1$ is randomisation within the pair. The
\emph{sensitivity value} $\Gamma^{*}$ is the smallest $\Gamma$ at which the worst-case p-value over
the model exceeds the level $\alpha$.
\end{definicao}

\prompt[conceito=gamma-de-rosenbaum,dificuldade=0.7]{State what $\Gamma$ bounds, why $\Gamma = 1$ is
the randomised experiment within pairs, and why the worst-case p-value can only grow with
$\Gamma$.}{$\Gamma$ bounds the ratio of the odds of treatment of two members of a pair that agree on
the observed covariates: a hidden confounder can tilt that chance by at most this factor. At
$\Gamma = 1$ both members have the same chance, which is a coin flip within the pair. The worst-case
p-value is a maximum over the admitted assignments, and the set admitted at $\Gamma'$ contains the
set admitted at any $\Gamma < \Gamma'$; a maximum over a larger set cannot decrease. So there is a
single point where significance is lost, $\Gamma^{*}$.}

\begin{antes}
\antesq[intervalo-conhecimento]{Which object of ours plays the worst-case p-value at each $\Gamma$,
and which plays $\Gamma^{*}$?}
\antesq[rval]{Which of our two radii has no counterpart in Rosenbaum?}
\end{antes}

\section{The translation, and where it breaks}

Inside a Markov equivalence class every graph generates the same observational distributions, and
Meek's rules certify that claims are consistent, not true: no test on the data separates an
orientation from its reverse. That is the position of Rosenbaum's unmeasured covariate, and it takes
the same answer, price the error instead of measuring it. The red rows are where the analogy breaks.

\begin{center}
\begin{tikzpicture}[
  lbox/.style={draw=laranjaesc, fill=laranjaesc!7, rounded corners=2pt, text width=5.2cm,
               minimum height=0.78cm, align=left, font=\scriptsize, inner sep=2.5pt},
  rbox/.style={draw=azulesc, fill=azulesc!6, rounded corners=2pt, text width=5.2cm,
               minimum height=0.78cm, align=left, font=\scriptsize, inner sep=2.5pt},
  xbox/.style={draw=vermesc, fill=vermesc!6, rounded corners=2pt, text width=5.2cm,
               minimum height=0.78cm, align=left, font=\scriptsize, inner sep=2.5pt},
  vbox/.style={draw=cinzamedio, dashed, rounded corners=2pt, text width=5.2cm,
               minimum height=0.78cm, align=center, font=\scriptsize, text=cinzamedio, inner sep=2.5pt},
  mid/.style={font=\scriptsize\sffamily\bfseries, text width=2.6cm, align=center, text=cinzaescuro},
  midx/.style={font=\scriptsize\sffamily\bfseries, text width=2.6cm, align=center, text=vermesc},
  y=0.81cm]
\node[font=\small\sffamily\bfseries, text=laranjaesc] at (-4.5,0.75) {Rosenbaum: matched pairs};
\node[font=\small\sffamily\bfseries, text=azulesc] at (4.5,0.75) {Our paper: orientation claims};
\foreach \yy/\mtxt/\ltxt/\rtxt/\rst/\mst in {
  0/{what the data\\ cannot test}/{an unmeasured covariate $u$ within pairs equal on the observed
     $x$}/{the orientations in $\Kset$ inside the class: Meek checks consistency, not
     truth}/rbox/mid,
  -1/{the parameter}/{$\Gamma \ge 1$: within-pair odds ratio of treatment}/{$r = 0,1,2,\dots$:
     the number of the analyst's claims allowed to be wrong}/rbox/mid,
  -2/{the report at\\ each value}/{the worst-case p-value over the assignments $\Gamma$
     admits}/{$\Ir{r}$: hull of the 95\,\% bands of the effects over the graphs of
     $\Ball{r}$}/rbox/mid,
  -3/{the breakdown\\ value}/{$\Gamma^{*}$: smallest $\Gamma$ at which significance is
     lost}/{$\rnul$: smallest $r$ at which the band contains zero}/rbox/mid,
  -4/{the scale}/{continuous: $\Gamma^{*} \approx 1.58$ is an ordinary report}/{an integer; on real
     structures $\rnul$ is almost always 1 or beyond 3}/xbox/midx,
  -5/{the reading}/{a parameter of a model for treatment assignment}/{a count with no probability
     reading: it does not estimate how many claims are wrong}/xbox/midx}
{
  \node[lbox] (L) at (-4.5,\yy) {\ltxt};
  \node[\mst] (M) at (0,\yy) {\mtxt};
  \node[\rst] (R) at (4.5,\yy) {\rtxt};
  \draw[-{Stealth[length=1.4mm]}, cinzamedio] (L.east) -- (M.west);
  \draw[-{Stealth[length=1.4mm]}, cinzamedio] (M.east) -- (R.west);
}
\node[vbox] at (-4.5,-6.15) {no counterpart};
\node[mid] at (0,-6.15) {a second radius};
\node[rbox] at (4.5,-6.15) {$\rval$: smallest $r$ at which the committed set $Z$ stops being
     valid};
\end{tikzpicture}
\end{center}

In symbols, with $\Gz = \mathrm{Meek}(\Chat, \Kset)$ and the committed set $Z = \Ostar(\Gz)$ read off
the analyst's graph, the right-hand column is this.

\begin{formulabox}[The translation, formally]
\[
\Ball{r} = \bigl\{\mathrm{Meek}(\Chat, \Kset \setminus S) : S \subseteq \Kset,\ |S| \le r\bigr\},
\qquad
\Ir{r} = \operatorname{hull} \bigcup_{G \in \Ball{r}} \mathrm{band}_{95}(G),
\]
\[
\rnul = \min\{r : 0 \in \Ir{r}\}, \qquad
\rval = \min\{r : Z \text{ is not valid in some } G \in \Ball{r}\},
\]
and each endpoint of $\Ir{r}$ comes with the set $S$ that produces it, named in the domain's
variables. Both events are upward closed in $r$ because $\Ball{r} \subseteq \Ball{r+1}$.
\end{formulabox}

A real analysis shows the translation end to end. The query G311 $\to$ G43 on \emph{magic-niab},
whose coefficients are fitted to real measurements, has two claims in the treatment's chain
component under the larger supplier's knowledge, and both are true.

\begin{medido}[One report, \texttt{results/interval/fig1\_candidates.json}, relaxed list (outside the registered rule), row 1]
True effect $-0.255$; estimate $-0.251$ $[-0.259, -0.244]$ at $n = 20\,000$. The band barely moves at
$r = 1$ (width ratio 1.06) and reaches zero at $r = 2$, when G1853 $\to$ G311 and G311 $\to$ G43 are
both retracted: $\rnul = 2$ on 20 of 20 seeds, $\rval = 2$. The supplier made six claims in the
network, three false, all outside the component.
\end{medido}

Read as Rosenbaum would: the conclusion that G311 lowers G43 rests on two named claims and survives the
withdrawal of either. Three breaks in the analogy have to be said once, in the paper. The scale is an
integer and, on real structures, almost without a middle (next section). There is no probability model
behind $r$: this row has $\rnul = 2$ with both local claims true and three false claims elsewhere, so
$r$ counts claims whose withdrawal changes the conclusion, not claims that are wrong. And the report is
honest only if the class contains the truth, which fails once the CPDAG is estimated
(Chapter~\ref{ch:pest}).

\prompt[conceito=intervalo-conhecimento,dificuldade=0.65]{Complete the translation: $\Gamma$
becomes what, the worst-case p-value becomes what, $\Gamma^{*}$ becomes what, and which object of
ours has no counterpart in Rosenbaum?}{$\Gamma$ becomes $r$, the number of the analyst's claims
allowed to be wrong. The worst-case p-value becomes the knowledge interval $\Ir{r}$, the hull of the
95\,\% bands over the graphs of the retraction ball $\Ball{r}$. $\Gamma^{*}$ becomes the null radius
$\rnul$, the smallest $r$ at which the band contains zero. The validity radius $\rval$ has no
counterpart: a matched design does not read an adjustment set off a graph that can break. The named
witnesses of each endpoint are also ours.}

\begin{antes}
\antesq[linhas-procedimento]{What do the rows of Table~1 have in common, and what do they differ
in?}
\antesq[colunas-cruzamento]{With one wrong claim, how much does the stated report cover, and at what
width does the one-claim interval restore coverage?}
\antesq[metrica-unica]{What does an $\Lnf$ above 1 mean?}
\end{antes}

\section{M1: Table 1 on real structure}

Before M1 the headline table existed only on synthetic generators, uncalibrated. Its
pre-registration (\texttt{results/interval/PREREGISTRATION.md}, P1 to P9) was written on 12 September
before \texttt{experiments/interval\_panel.py} produced a row. A run takes about 90 seconds; five runs
with 3 to 7 processes gave a byte-identical \texttt{rows.csv}.

\textbf{Design.} The frozen frame, 528 queries. Panel~A: the 23 other networks of the ledger, real
structure, semi-synthetic coefficients (magnitude uniform on $[0.5, 1.5]$, random sign, unit noise),
448 queries. Panel~B: \texttt{ecoli70}, \texttt{arth150}, \texttt{magic-niab}, \texttt{magic-irri}
with fitted coefficients, 80 queries. Knowledge: controlled reversals of $b = 0, 1, 2$ true claims
inside the treatment's chain component (a reversal outside it moves no interval, a reduction tested
against brute force on 7,770 pairs), and the 7B and 72B suppliers' elicited claims stratified by false
local claims. Samples at $n = 1000$ and $20\,000$, twenty seeds. \textbf{One estimator} in every row,
the OLS coefficient of $X$ given a set; rows differ only in the graphs the hull is taken over: the
committed report (PLAIN), a matched control retracting one claim far from $X$ and $Y$ (FAR1),
$\Ir{1}$, $\Ir{2}$, and the whole class (BLANKET, discarding the knowledge). The metric $\Lnf$
stretches each procedure's intervals by the smallest common factor giving 95\,\% coverage in the cell
and divides the mean stretched width by BLANKET's (Chapter~\ref{ch:tables} opens a cell).
\textbf{Above 1 means wider than throwing the knowledge away.}

\begin{table}[!ht]
\centering\footnotesize
\setlength{\tabcolsep}{3.4pt}
\caption{M1, informative queries, $n = 1000$: $\Lnf$ [network-bootstrap 95\,\% interval] / raw
coverage. The three columns are different populations (see the rows line). Source:
\texttt{results/interval/TABLE1\_L95.md}.}
\label{tab:m1-gamma}
\begin{tabular}{@{}l ll ll ll@{}}
\toprule
 & \multicolumn{2}{c}{$b = 0$} & \multicolumn{2}{c}{$b = 1$} & \multicolumn{2}{c}{$b = 2$} \\
\cmidrule(lr){2-3}\cmidrule(lr){4-5}\cmidrule(lr){6-7}
procedure & $\Lnf$ & cov. & $\Lnf$ & cov. & $\Lnf$ & cov. \\
\midrule
\multicolumn{7}{@{}l}{\emph{Panel~A, semi-synthetic coefficients, 23 networks}} \\
PLAIN (committed) & 0.464 [0.35, 0.57] & 0.967 & 5.123 [3.24, inf) & 0.597 & 4.082 [2.71, inf) & 0.554 \\
FAR1 (control)   & 0.655 & 0.967 & 10.771 & 0.640 & 7.539 & 0.614 \\
$\Ir{1}$         & 0.957 [0.87, 0.99] & 0.982 & 0.851 [0.76, 0.94] & 0.980 & 4.234 & 0.781 \\
$\Ir{2}$         & 0.984 & 0.984 & 0.980 & 0.984 & 0.847 [0.74, 0.94] & 0.984 \\
BLANKET (class)  & 1.000 & 0.984 & 1.000 & 0.985 & 1.000 & 0.988 \\
rows (networks)  & 391 (23) & & 359 (23) & & 131 (9) & \\
\midrule
\multicolumn{7}{@{}l}{\emph{Panel~B, fitted coefficients, 4 networks}} \\
PLAIN            & 0.840 [0.76, 0.88] & 0.969 & 3.316 [0.81, 4.59] & 0.885 & 2.588 & 0.891 \\
$\Ir{1}$         & 0.979 & 0.979 & 0.974 [0.94, 1.00] & 0.980 & 0.887 & 0.947 \\
$\Ir{2}$         & 0.997 & 0.980 & 0.989 & 0.983 & 0.945 & 0.979 \\
rows (networks)  & 69 (4) & & 44 (4) & & 35 (4) & \\
\bottomrule
\end{tabular}
\end{table}

On Panel~A, with no wrong claim the committed report is the narrowest (0.464) and the one-claim
interval pays a premium (0.957): insurance costs when nothing is wrong (P1). With one wrong claim the
committed report covers 0.597 and would need 5.1 class widths to reach 95\,\%, while $\Ir{1}$ covers
0.980 at 0.851 of the class width. With two, $\Ir{1}$ is not enough (0.781) and $\Ir{2}$ pays.
\figref{fig:l95} draws the crossing.

\begin{figure}[!ht]
\centering
\includegraphics[width=0.74\linewidth]{l95_crossing.pdf}
\caption{Panel~A: with one wrong claim the committed report costs more than five class widths and
retracting a far claim does not help; with two, only the budget-two procedures stay below 1. Panel~B's
four networks keep the order at smaller magnitudes. The columns are different populations, so the lines
between them do not follow the same analyses.}
\label{fig:l95}
\end{figure}

\begin{medido}[The crossing and the saving, \texttt{results/interval/TABLE1\_L95.md} and \texttt{PREREGISTRATION.md} A.2]
Against discarding the knowledge, $\Ir{1}$ at $b = 1$ saves 14.9\,\% of the width at $n = 1000$ and
24.2\,\% at $n = 20\,000$ on Panel~A (0.851 and 0.758), and 2.6\,\% and 4.8\,\% on Panel~B (0.974 and
0.952). With errors nobody planted (elicited, one false local claim, suppliers pooled, $n = 1000$)
the order repeats: PLAIN 0.600 at $\Lnf$ 6.92 against 0.867 for $\Ir{1}$ on Panel~A (208 rows), 0.696
at 2.41 against 0.901 on Panel~B (23 rows). P1, P2, P4, P5, P7, P8, P9 supported; P3 partial; P6
falsified.
\end{medido}

The saving is modest, and the paper should say so: the interval covers where the committed report does
not, at less than the price of discarding the knowledge.

\prompt[conceito=colunas-cruzamento,dificuldade=0.75]{Panel~A, $n = 1000$, one wrong claim: give the
coverage and $\Lnf$ of the committed report and of $\Ir{1}$, and say what 5.12 means.}{PLAIN covers
0.597 with $\Lnf = 5.123$; $\Ir{1}$ covers 0.980 with $\Lnf = 0.851$ [0.76, 0.94], over 359
informative queries on 23 networks. The 5.12 means that to reach 95\,\% coverage the committed
report's intervals would have to be stretched to a mean width of 5.1 times the calibrated width of the
whole class: trusting one wrong claim costs more than discarding all the knowledge.}

\begin{antes}
\antesq[r0-binario]{Where an effect is established with the full knowledge, what values does
$\rnul$ take on real structures?}
\antesq[banda-plugin]{How often does the hull of point estimates state a null radius above the
population one, and how often does the band?}
\end{antes}

\section{What M1 changed in the report}

The translation promises a scale: how many wrong claims until the conclusion falls. On real
structures the scale has almost no middle.

\begin{medido}[$\rnul$ by knowledge condition, counted in \texttt{results/interval/rows.csv} (\texttt{r0\_pop})]
Frame queries with an effect established under the full knowledge ($\rnul \ge 1$); ``beyond 3''
pools \texttt{>3}, \texttt{never} and censored. Panel~A, true claims ($b = 0$): 221 queries, $\rnul = 1$
on 71, 2 or 3 on 16, beyond 3 on 134, so 1 or beyond 3 on 92.8\,\%. The same share is 88.8\,\% at
$b = 1$, 74.2\,\% at $b = 2$ (66 queries), 93.6\,\% (7B), 94.1\,\% (72B) and 97.7\,\% under the
recovering set (6 of 257 at 2 or 3). Panel~B: 25 of 25 at $b = 0$, 33 of 38 under 72B. Restricted to
informative queries the $b = 0$ share on Panel~A is 148 of 164 (90.2\,\%): a non-informative query
with a nonzero effect is ``never'' by construction.
\end{medido}

\begin{figure}[!ht]
\centering
\includegraphics[width=0.70\linewidth]{r0_distribution.pdf}
\caption{Under every knowledge source and on both panels, where an effect is established the
conclusion rests on one claim or on none within depth three: the mass at 2 and 3 is a sliver, and
even complete oracle knowledge does not spread it.}
\label{fig:r0}
\end{figure}

\figref{fig:r0} shows it as a distribution. In the analyst's words, \emph{the conclusion rests on no
claim, or on one claim that can be named}. That changed decision D2 in the outline of 13 September:
lead with the knowledge interval and Table~1, print $\rnul$ beside it as ``rests on no claim / on this
claim'' with its distribution, and retire the title candidate that named the null radius.

A reading trap: $\rnul = 0$ (227 of 391 informative Panel~A queries under true claims, in
\texttt{summary.json}) means there is no established effect to break, not a break at the first
retraction.

\textbf{The band, not the hull of point estimates, is the estimator} (P7). At $n = 1000$, over 23,120
replicates of the controlled conditions, the plug-in hull states a null radius above the population
one on 39.7\,\% of replicates and the band on 1.7\,\% (39.0\,\% and 1.8\,\% at $n = 20\,000$). A
possible effect near zero lands on one side of it by sampling error, and the plug-in errs towards
claiming robustness.

\textbf{A confounding bound at its own oracle} (M2, P9). The Cinelli and Hazlett bound, given the true
partial $R^{2}$ of the block missing from $Z$, is defined on 666 of 808 committed-set queries (on 101
the wrong claim put a true descendant into $Z$; on 41 no valid set exists). There it is narrower than
$\Ir{1}$ (Panel~A, $b = 1$, $n = 20\,000$: $\Lnf$ 0.445 against 0.632), being handed what the analyst
lacks. On 102 queries whose set holds a true descendant and is invalid it covers 13.4\,\% at
$n = 20\,000$: a wrong orientation is outside its model class.

\begin{fonte}[Rosenbaum 2002; Cinelli and Hazlett 2020; Masten and Poirier 2020]
Rosenbaum, \emph{Observational Studies}, 2nd ed.: the chapter on hidden bias, where our report's form
comes from. Cinelli and Hazlett, JRSS~B 82(1):39--67: the bound M2 implemented. Masten and Poirier's
breakdown frontiers, \emph{Quantitative Economics} 11(1):41--111: a breakdown value with inference,
outside graphs, not yet in your bibliography.
\end{fonte}

\textbf{Table~2 came out thinner than planned.} Seven of the nine declared queries have no claim in the
treatment's component, so $\rnul$ is ``never'' (\texttt{results/interval/TABLE2\_REPORTS.md}); only
\texttt{mediator} has a finite one (2). The informative reports are the 17 real-coefficient queries,
mostly $\rnul = 1$ with the claim named (Chapter~\ref{ch:tables}).


\begin{sonorio}[The analogy stated once, with the place it breaks]
Following Rosenbaum, we report how many of the analyst's orientation claims would have to be withdrawn
before the conclusion changes, beside the interval of effects those withdrawals leave possible. Unlike
a sensitivity value, this count is an integer with no probability reading; on the real structures
measured here it is almost always one or more than three, so the report names the claim on which the
conclusion rests.
\end{sonorio}

\prompt[conceito=r0-binario,dificuldade=0.7]{Give the measured count that shows $\rnul$ near-binary,
with its population, and say what it changed in the way the paper reports.}{Panel~A, true claims
($b = 0$), frame queries with an established effect: 205 of 221 (92.8\,\%) have $\rnul = 1$ (71) or
beyond 3 (134); restricted to informative queries 148 of 164 (90.2\,\%); Panel~B 25 of 25. So the paper
leads with the knowledge interval and Table~1 and reports $\rnul$ as ``the conclusion rests on no
claim, or on this named claim'', printing the distribution; the title that named the radius was
withdrawn.}

\naopasse[nulidade]{Name the step where ``significance is lost'' becomes ``the band contains zero'', and
say why the band and not the plug-in hull, with the P7 numbers and their sample size.}

\subsection*{Exercises}

\begin{exercicio}[metodo=calculation,conceito=intervalo-conhecimento]
Illustrative, population level. Local claims $\Kset = \{a, b, c\}$, committed effect 0.80; effects
each retraction leaves possible: $\emptyset$: $\{0.80\}$; $\{a\}$: $\{0.80\}$; $\{b\}$: $\{0.80, 0.35\}$; $\{c\}$:
$\{0.80\}$; $\{a,b\}$: $\{0.80, 0.35\}$; $\{a,c\}$: $\{0.80, 0.00\}$; $\{b,c\}$:
$\{0.80, 0.35, 0.00\}$; $\{a,b,c\}$ (the class): $\{0.80, 0.35, 0.00, -0.10\}$. $Z$ is invalid after
retracting $b$ and valid after retracting $a$ or $c$ alone. (a) Compute $\Ir{0}$ to $\Ir{3}$. (b)
Give $\rnul$ and $\rval$, with the witnesses. (c) What share of the class width does $\Ir{1}$ cover?
\end{exercicio}
\begin{solucao}
(a) $\Ir{0} = [0.80, 0.80]$; $\Ir{1} = [0.35, 0.80]$ (from $\{b\}$); $\Ir{2} = [0.00, 0.80]$ (from
$\{a,c\}$ or $\{b,c\}$); $\Ir{3} = [-0.10, 0.80]$, the class. (b) $0 \notin \Ir{1}$ and $0 \in \Ir{2}$,
so $\rnul = 2$, witnessed by $\{a,c\}$ (and $\{b,c\}$); no single claim reaches zero, but $a$ and $c$
together do, an interaction. $\rval = 1$, witnessed by $b$. Here the set breaks before the sign does,
as on \texttt{mediator} ($\rval = 1$, $\rnul = 2$). (c) $0.45 / 0.90 = 50\,\%$. The listed sets respect
the pre-registered reduction ``retracting more never removes an extension'' (each superset retraction
contains its subsets' effects), which is what makes $\Ir{r}$ non-decreasing.
\end{solucao}

\begin{exercicio}[metodo=drawing,conceito=testemunha]
Draw the report for \texttt{magic-niab}, G311 $\to$ G43 ($n = 20\,000$, seed 0), effect against
$r = 0, 1, 2$, under two claim sets: (i) the 72B supplier's two local claims,
bands $[-0.260, -0.245]$, $[-0.260, -0.244]$, $[-0.260, 0.000]$; (ii) the recovering set, whose only
local claim is G1853 $\to$ G311, bands $[-0.260, -0.245]$ at $r = 0$ and $[-0.260, 0.000]$ at $r = 1$
and $r = 2$ (\texttt{results/interval/profiles.jsonl}). Mark $\rnul$ and its witness, and explain how one query has $\rnul = 2$ in one report and 1 in the
other.
\end{exercicio}
\begin{solucao}
\begin{center}
\begin{tikzpicture}[x=1.7cm, y=5.2cm]
\foreach \dx/\a/\b/\c/\col in {-0.1/-0.245/-0.244/0.0/azulesc, 0.1/-0.245/0.0/0.0/laranjaesc}{
  \draw[\col, line width=2.4pt] (0+\dx,-0.260) -- (0+\dx,\a);
  \draw[\col, line width=2.4pt] (1+\dx,-0.260) -- (1+\dx,\b);
  \draw[\col, line width=2.4pt] (2+\dx,-0.260) -- (2+\dx,\c);
}
\draw[cinzaescuro] (-0.5,0) node[left, font=\scriptsize] {zero} -- (2.5,0);
\draw[cinzamedio, densely dotted] (-0.5,-0.255) node[left, font=\scriptsize] {truth $-0.255$} -- (2.5,-0.255);
\foreach \r in {0,1,2} \node[font=\scriptsize] at (\r,-0.30) {$r = \r$};
\node[font=\scriptsize, text=laranjaesc, align=right, anchor=east] at (0.9,-0.10)
  {(ii) recovering set: $\rnul = 1$\\ witness G1853$\to$G311};
\node[font=\scriptsize, text=azulesc, align=left, anchor=west] at (2.2,-0.10)
  {(i) 72B supplier: $\rnul = 2$\\ witness G1853$\to$G311\\ and G311$\to$G43};
\end{tikzpicture}
\end{center}
Both start from a point interval at $r = 0$, so both analyst graphs orient G311 $\to$ G43. In (ii) that
orientation is not asserted, so it is propagated from G1853 $\to$ G311, and retracting that one claim
admits the graph where G43 is a parent of G311 and the effect is zero. In (i) it is also asserted
directly, so it survives either single retraction and falls only when both go. $\rnul$ belongs to the
claims around the query, not to the query, so a Figure~1 caption must name the claim set.
\end{solucao}

\section*{Chapter map}
\begin{center}
\begin{tikzpicture}[node distance=0.6cm and 1.15cm]
\node[cmapc] (g) {$\Gamma$: how strong\\ hidden bias must be};
\node[cmap, right=of g] (k) {orientation claims:\\ untestable in the class};
\node[cmapc, right=of k] (t) {$r$, $\Ir{r}$, $\rnul$, $\rval$,\\ named witnesses};
\node[cmapg, below=of t] (m1) {M1: $\Ir{1}$ covers 0.980\\ at 0.851 of the class};
\node[cmapr, below=of k] (bin) {$\rnul$ near-binary:\\ 1 or beyond 3};
\node[cmap, below=of g] (d2) {D2: interval first,\\ $\rnul$ as the named claim};
\node[cmap, below=of bin] (ch) {oracle $R^{2}$ bound:\\ 13.4\,\% on descendants};
\draw[cedge] (g) -- node[above]{same move} (k);
\draw[cedge] (k) -- node[above]{translated into} (t);
\draw[cedge] (t) -- node[right]{measured by} (m1);
\draw[cedge] (m1) -- node[above]{shows} (bin);
\draw[cedge] (bin) -- node[above]{decides} (d2);
\draw[cedge] (ch) -- node[right]{other model class} (bin);
\end{tikzpicture}
\end{center}
